\section{Luce in forma particellare}
Come abbiamo appena visto, se si suppone che la luce è in forma esclusivamente ondulatoria si va in contro ad alcuni conflitti con i dati sperimentali. È quindi necessario ri-definire il concetto stesso di "onda luminosa": ora non sarà più una semplice onda ma un pacchetto d'onda formato da una distribuzione continua di onde monocromatiche. È facile immaginare che la velocità del pacchetto d'onda sarà la velocità della luce: prima avevamo che per un'onda monocromatica $f(x,t) = a \cos (k \cdot x - \omega t)$ la velocità di propagazione era data da $v = |k|/\omega = c$ nel caso della luce. In particolare si aveva $\omega = ck$, che viene ora chiamata \textbf{relazione (o legge) di disperzione} \index{legge di disperzione} del pacchetto d'onda.\index{Relazione di dispersione} In generale si nota che vale $v = \deriv[k]{\omega(k)}$; questa quantità viene chiamata \textbf{velocità di gruppo}.\index{Velocità di Grupo}

\subsection*{Distribuzione particellare}
De Broglie: data una distribuzione di particelle, le cui cartteristiche seguono quindi le caratteristiche della meccanica classica, si ha che una particella è definita da \textbf{massa} $m$, \textbf{quantità di moto} $\vec p$, \textbf{energia} $E$.
Essa ha anche un'onda associata, come visto prima, e in particolare ogni cosa ha un'onda associata! possibilmente stazionaria, anche la materia:
\[
	\omega = \frac{E}{\hslash}, \ \ \ \ 
	\vec k = \frac{\vec p}{\hslash}
\]
Infatti si può fare l'esperimento di Young con fasci di materia e produrre anche in questo caso delle figure di diffrazione! Ovviamente sono necessarie condizioni molto particolari, in quanto le lunghezze d'onda degli oggetti materiali sono molto più piccole di quelle della luce.

\subsection*{Legge di dispersione delle Matter Waves}
Per la luce era \( \omega = c|\vec k| \)\\
Riarrangiando le espressioni introdotte poco fa risulta
\[
	\omega = \frac{\hslash|\vec k|^2}{2m}
\]
che è una legge di dispersione quadratica!

Pacchetti d'onda? Facile!
risulta
\[
	\vec v_g = \Grad _{\vec k} \omega = \frac{2 \hslash \vec k}{2 m} = \frac{\vec p}{m} = \vec v
\]

che è proprio la velocità della particella. Al chè
\begin{align*}
	\psi(\vec x, t) = e^{i(\vec k \scal \vec x - \omega t)} \\
	\tderiv{\psi} = -i\omega \psi = \frac{E}{\hslash}\psi \\
	\Lap \psi = -|k|^2\psi
\end{align*}
E quindi risulta (\textbf{Equazione di Schroedinger}) \index{Equazione di Schrödinger}
\begin{equation}
	i\hslash \tderiv{\psi (\vec x, t)} = - \frac{\hslash^2}{2m} \Lap \psi(\vec x, t)
\end{equation}
La soluzione a quest'equazione differenziale rappresenta lo \textbf{stato} del sistema, ovvero una quantità che permette di derivare la probabilità di trovare la particella in una certa posizione nello spazio ad un certo istante di tempo.
In sostanza stiamo affermando che le misurazioni effettuabili su una matter wave sono regolate da principi statistici! \\
È la quantità $|\psi(\vec x, t)|^2 $ ad essere una vera e propria distribuzione di probabilità, quindi si pone
\[
	\int_{\R^3} |\psi(\vec x, t)|^2 dx = 1 = ||\psi||^2
\]
dove $||\scal||$ è la norma $L^2$
Risulta comodo definire anche l'operatore Hamiltoniano \index{operatore Hamiltoniano}
\begin{align*}
	\H = &- \frac{\hslash^2}{2m} \Lap + U(\vec x)\\
	\H\psi = - \frac{\hslash^2}{2m} &\Lap \psi(\vec x, t) + U(\vec x) \psi(\vx, t)\\
	\vec p = &- i \hslash \Div
\end{align*}
per cui, dato il prodotto scalare tipico di $L^2(\C)$
\[
	\fscal f g = \int_{\R^3} f^\star g
\]
bisogna dimostrare (esercizio) che $\fscal\phi {\H\phi}$ = $\fscal {\H\phi} \phi \in \R$, ovvero è un operatore \textbf{hermitiano}.\\
La funzione $U(\vec x)$ è un potenziale definito localmente, utile per poter introdurre dei vincoli sul moto del pacchetto d'onda.

La funzione $\psi$ è chiamata \textbf{funzione d'onda} \index{Funzione d'onda}, ed è quell'oggetto che contiene tutte le informazioni sullo stato del sistema (per convenienza viene anche chiamata \textit{stato}).

Per finire, siccome $\sfabs \psi$ è una distribuzione di probabilità volumetrica, è opportuno definire un "flusso di probabilità", ovvero una funzione che esprime localmente la variazione nel tempo della densità di probabilità. Per esercizio si riccavi la formula
\[
	\J = \frac{\hslash}{2im}[\psi^*(\Div\psi) - \psi(\Div\psi^*)]
\]
Si ricavi inoltre che questa quantità coincide con la velocità della particella.

L'ultimo fatto interessante è che le matter waves possono interferire tra loro.